% ============================================================
%  CS7IS2 — Artificial Intelligence — Assignment 3
%  Minimax & Reinforcement Learning for Board Games
%  Trinity College Dublin  ·  2025–2026
% ============================================================
\documentclass[12pt, a4paper]{article}

\usepackage[top=2.5cm, bottom=2.5cm, left=2.8cm, right=2.8cm]{geometry}
\usepackage{graphicx}
\graphicspath{{../results/figures/}{screenshots/}{../results/figures/screenshots/}{./}}
\usepackage{booktabs}
\usepackage{array}
\usepackage{multirow}
\usepackage{amsmath}
\usepackage{amssymb}
\usepackage{xcolor}
\usepackage{hyperref}
\usepackage{fancyhdr}
\usepackage{setspace}
\usepackage{caption}
\usepackage{subcaption}
\usepackage{enumitem}
\usepackage{titlesec}
\usepackage{parskip}
\usepackage{microtype}
\usepackage{float}
\usepackage{listings}
\usepackage{algorithmicx}
\usepackage{algpseudocode}
\usepackage{algorithm}
\usepackage[super,sort&compress]{natbib}
\usepackage{tocloft}

\renewcommand{\cftsecleader}{\cftdotfill{\cftnodots}}
\renewcommand{\cftsubsecleader}{\cftdotfill{\cftnodots}}
\renewcommand{\cftsecfont}{\bfseries\color{tcdblue}}
\renewcommand{\cftsecpagefont}{\bfseries\color{tcdblue}}
\setlength{\cftsecnumwidth}{2.5em}
\setlength{\cftsubsecnumwidth}{3.2em}
\setlength{\cftsubsubsecnumwidth}{4.0em}
\setlength{\cftbeforesecskip}{6pt}
\setlength{\cftbeforesubsecskip}{2pt}

\definecolor{tcdblue}{RGB}{0, 45, 114}
\definecolor{tcdgold}{RGB}{163, 127, 58}
\definecolor{codegrey}{RGB}{240, 240, 240}
\definecolor{linkblue}{RGB}{0, 85, 160}

\hypersetup{colorlinks=true, linkcolor=linkblue, citecolor=tcdblue, urlcolor=linkblue,
    pdftitle={CS7IS2 Assignment 3}, pdfauthor={Naman Kumar}}

\titleformat{\section}{\large\bfseries\color{tcdblue}}{{\thesection}}{0.8em}{}[\vspace{-4pt}\rule{\linewidth}{0.6pt}]
\titleformat{\subsection}{\normalsize\bfseries\color{tcdblue}}{\thesubsection}{0.8em}{}
\titleformat{\subsubsection}{\normalsize\bfseries}{\thesubsubsection}{0.8em}{}

\pagestyle{fancy}\fancyhf{}
\fancyfoot[L]{%
    \IfFileExists{trinity_logo.png}{%
        \includegraphics[height=0.75cm]{trinity_logo.png}%
    }{\IfFileExists{trinity_logo.jpg}{%
        \includegraphics[height=0.75cm]{trinity_logo.jpg}%
    }{{\footnotesize\color{tcdblue}\textbf{TCD}}}}%
}
\fancyfoot[C]{\small\color{tcdblue} CS7IS2 — AI \textbf{·} Assignment 3}
\fancyfoot[R]{\small\color{tcdblue}\thepage}
\renewcommand{\headrulewidth}{0pt}
\renewcommand{\footrulewidth}{0.4pt}
\setlength{\headheight}{15pt}
\setlength{\footskip}{30pt}

\lstset{backgroundcolor=\color{codegrey}, basicstyle=\ttfamily\small,
    keywordstyle=\bfseries\color{tcdblue}, commentstyle=\itshape\color{gray},
    stringstyle=\color{tcdgold}, breaklines=true, frame=single, numbers=left,
    numberstyle=\tiny\color{gray}, tabsize=4, captionpos=b}

\captionsetup{font=small, labelfont={bf,color=tcdblue}, margin=10pt}

% ── flexible image loader: tries .png then .jpg then shows placeholder ───────
\newcommand{\tryimg}[3][width=\linewidth]{%
    \IfFileExists{#2.png}{%
        \includegraphics[#1]{#2.png}%
    }{%
        \IfFileExists{#2.jpg}{%
            \includegraphics[#1]{#2.jpg}%
        }{%
            \IfFileExists{#2.jpeg}{%
                \includegraphics[#1]{#2.jpeg}%
            }{%
                \fbox{\parbox{0.6\linewidth}{\centering\vspace{1.8cm}%
                    \textit{[Missing image: #2.png / .jpg]}%
                \vspace{1.8cm}}}%
            }%
        }%
    }%
}
\newcolumntype{C}[1]{>{\centering\arraybackslash}p{#1}}
\setlength{\parskip}{6pt}
\setlist[itemize]{noitemsep, topsep=2pt}
\setlist[enumerate]{noitemsep, topsep=2pt}


\begin{document}
\begin{titlepage}\centering\vspace*{0.5cm}
    \IfFileExists{trinity_logo.png}{%
        \includegraphics[width=5cm]{trinity_logo.png}\\[0.8cm]
    }{\IfFileExists{trinity_logo.jpg}{%
        \includegraphics[width=5cm]{trinity_logo.jpg}\\[0.8cm]
    }{%
        {\fontsize{28}{32}\selectfont\color{tcdblue}\textbf{Trinity College Dublin}}\\[0.4cm]
        {\large\color{tcdgold} The University of Dublin}\\[0.8cm]
        \rule{0.6\linewidth}{1.5pt}\\[0.8cm]
    }}
    {\Large\bfseries\color{tcdblue} CS7IS2 — Artificial Intelligence}\\[0.3cm]
    {\large Assignment 3: Minimax and Reinforcement Learning}\\[1.2cm]
    \rule{\linewidth}{0.6pt}\\[0.6cm]
    {\LARGE\bfseries\color{tcdblue} Comparative Analysis of Minimax}\\[0.3cm]
    {\LARGE\bfseries\color{tcdblue} and Reinforcement Learning}\\[0.3cm]
    {\large\bfseries\color{tcdblue} on Tic Tac Toe and Connect 4}\\[0.6cm]
    \rule{\linewidth}{0.6pt}\\[1.2cm]
    \begin{tabular}{rl}
        \textbf{Author}      & Naman Kumar\\[4pt]
        \textbf{Student ID}  & 25340053\\[4pt]
        \textbf{Module}      & CS7IS2 — Artificial Intelligence\\[4pt]
        \textbf{Institution} & Trinity College Dublin\\[4pt]
        \textbf{Date}        & \today\\
    \end{tabular}
    \vfill
\end{titlepage}

\tableofcontents
\newpage


% ══════════════════════════════════════════════════════════════════════════════
\section{Introduction}

Game playing has been a central challenge in artificial intelligence research since the field's
inception. Two-player, zero-sum, perfect-information games provide a controlled setting in which
to study decision-making under adversarial conditions, where every gain for one player is a
corresponding loss for the other~\cite{vonneumann1944}.

This report implements and compares two fundamentally different paradigms for playing board games.
The first paradigm, \textbf{adversarial search}, constructs a game tree and evaluates positions by
simulating all possible future moves. The canonical algorithm is Minimax~\cite{russell2010},
further improved by alpha-beta pruning~\cite{knuth1975} to reduce the number of nodes evaluated.
The second paradigm, \textbf{reinforcement learning (RL)}, learns a policy by repeatedly playing
the game and updating value estimates based on outcomes. Two RL approaches are considered:
tabular Q-learning~\cite{watkins1992}, which stores state-action values in an explicit lookup
table, and Deep Q-Networks (DQN)~\cite{mnih2015}, which use a neural network to approximate
the Q-function.

These algorithms are evaluated on two games of very different complexity.
\textbf{Tic Tac Toe} (TTT) is a $3\times3$ grid game with approximately 5,478 legal positions
and a maximum of 9 moves, making full-depth search entirely feasible.
\textbf{Connect 4} (C4) is a $6\times7$ drop-piece game with approximately
$4.5 \times 10^{12}$ positions~\cite{allis1988} and a maximum of 42 moves, rendering full-depth
minimax computationally infeasible on consumer hardware.

The report is structured as follows: Section~2 describes the game environments and default
opponent. Section~3 details each algorithm. Section~4 specifies the experimental setup.
Section~5 analyses Connect~4 scalability. Sections~6 and~7 present results for TTT and C4.
Section~8 reports parameter sensitivity experiments. Section~9 provides an overall comparison,
and Section~10 states the conclusions.


% ══════════════════════════════════════════════════════════════════════════════
\section{Game Environments}

\subsection{Tic Tac Toe}

Tic Tac Toe is played on a $3\times3$ grid. Two players alternate placing their marks (X or O);
the first to align three marks in a row, column, or diagonal wins. If all nine cells are filled
without a winner, the game is a draw.

The game has an upper bound of $9! = 362{,}880$ possible move sequences, but many are equivalent
or terminate early. The number of distinct legal board positions is approximately
5,478~\cite{russell2010}. The branching factor starts at 9 and decreases by 1 per ply, giving an
average branching factor of roughly 4. The game tree is shallow enough that Minimax can
exhaustively evaluate every outcome within milliseconds.

\begin{figure}[H]
    \centering
    \tryimg[width=0.45\linewidth]{screenshots/ttt_board}
    \caption{Tic Tac Toe game in progress. X has played three moves; O has played two.}
    \label{fig:ttt_board}
\end{figure}


\subsection{Connect 4}

Connect 4 is played on a vertical $6\times7$ board. Players alternate dropping coloured discs
into one of seven columns; each disc falls to the lowest available row. The first player to
connect four discs horizontally, vertically, or diagonally wins. If all 42 cells are filled
without a winner, the game is a draw.

The game's complexity dwarfs Tic Tac Toe: the number of legal positions is approximately
$4.5 \times 10^{12}$~\cite{allis1988}, and the branching factor is up to 7 throughout most of
the game. With a maximum game depth of 42, the game tree is vastly too large for exhaustive
search. Allis proved that the first player can force a win with perfect play, but the proof
required months of computation~\cite{allis1988}. This complexity difference motivates the use
of depth-limited minimax and RL approaches that generalise from limited training.

\begin{figure}[H]
    \centering
    \tryimg[width=0.55\linewidth]{screenshots/c4_board}
    \caption{Connect 4 game in progress showing a typical mid-game position.}
    \label{fig:c4_board}
\end{figure}


\subsection{Default Opponent}
\label{sec:default_opp}

The assignment requires an opponent better than fully random. The default opponent implements
a three-tier decision rule:

\begin{enumerate}
    \item \textbf{Win}: if a move exists that immediately wins, take it.
    \item \textbf{Block}: if the adversary has a winning move next turn, block it.
    \item \textbf{Positional preference}: otherwise, prefer the centre cell (TTT) or centre
          columns (C4), as these offer the most strategic flexibility.
\end{enumerate}

This opponent never misses a one-move win or allows a one-move loss, making it meaningfully
stronger than random while still being beatable by competent algorithms.

\begin{algorithm}[H]
\caption{Default Opponent Strategy}
\begin{algorithmic}[1]
\Function{ChooseMove}{game}
    \For{each valid move $m$}
        \If{$m$ wins the game for me} \Return $m$ \EndIf
    \EndFor
    \For{each valid move $m$}
        \If{$m$ would win the game for opponent} \Return $m$ \Comment{Block} \EndIf
    \EndFor
    \State \Return move closest to centre
\EndFunction
\end{algorithmic}
\end{algorithm}


% ══════════════════════════════════════════════════════════════════════════════
\section{Algorithm Implementations}

\subsection{Minimax (Without Pruning)}

Minimax is an adversarial search algorithm that models a two-player game as a tree where
players alternate between maximising and minimising the utility
value~\cite{russell2010,vonneumann1944}. Starting from the current state, the algorithm
recursively explores every possible move sequence until terminal states are reached. Terminal
states are scored relative to the maximising player: a win returns a positive value, a loss a
negative value, and a draw returns zero. The score is adjusted by depth to prefer faster wins
and slower losses (e.g., $10 - \text{depth}$ for TTT, $1000 - \text{depth}$ for C4).

At each non-terminal node, the maximising player selects the child with the highest value,
while the minimising player selects the lowest. This guarantees optimal play against a perfect
adversary.

\begin{algorithm}[H]
\caption{Minimax (No Pruning)}
\begin{algorithmic}[1]
\Function{Minimax}{game, depth, isMaximizing}
    \If{game is terminal}
        \State \Return utility(game, depth)
    \EndIf
    \If{isMaximizing}
        \State $\text{best} \leftarrow -\infty$
        \For{each valid move $m$}
            \State apply $m$; \;
            $\text{best} \leftarrow \max(\text{best},\;\Call{Minimax}{game, depth{+}1, \text{False}})$;
            \; undo $m$
        \EndFor
    \Else
        \State $\text{best} \leftarrow +\infty$
        \For{each valid move $m$}
            \State apply $m$; \;
            $\text{best} \leftarrow \min(\text{best},\;\Call{Minimax}{game, depth{+}1, \text{True}})$;
            \; undo $m$
        \EndFor
    \EndIf
    \State \Return best
\EndFunction
\end{algorithmic}
\end{algorithm}


\subsection{Minimax with Alpha-Beta Pruning}

Alpha-beta pruning optimises Minimax by maintaining two bounds: $\alpha$ (the best value the
maximiser can guarantee) and $\beta$ (the best value the minimiser can guarantee). When
$\beta \leq \alpha$ at any node, remaining children cannot influence the result and are
pruned~\cite{knuth1975,russell2010}.

Alpha-beta never changes the result — it returns the same optimal move — but can dramatically
reduce nodes evaluated. In the best case (perfect move ordering), it reduces the effective
branching factor from $b$ to $\sqrt{b}$, squaring the searchable depth in the same
time~\cite{knuth1975}. Even with random ordering, significant savings are typical.

\begin{algorithm}[H]
\caption{Minimax with Alpha-Beta Pruning}
\begin{algorithmic}[1]
\Function{AlphaBeta}{game, depth, $\alpha$, $\beta$, isMax}
    \If{terminal or depth $\geq$ limit}
        \State \Return utility or evaluate(game)
    \EndIf
    \If{isMax}
        \State $v \leftarrow -\infty$
        \For{each valid move $m$}
            \State apply $m$; \;
            $v \leftarrow \max(v,\;\Call{AlphaBeta}{game, depth{+}1, \alpha, \beta, \text{False}})$;
            \; undo $m$
            \State $\alpha \leftarrow \max(\alpha, v)$
            \If{$\beta \leq \alpha$} \textbf{break} \EndIf
        \EndFor
    \Else
        \State $v \leftarrow +\infty$
        \For{each valid move $m$}
            \State apply $m$; \;
            $v \leftarrow \min(v,\;\Call{AlphaBeta}{game, depth{+}1, \alpha, \beta, \text{True}})$;
            \; undo $m$
            \State $\beta \leftarrow \min(\beta, v)$
            \If{$\beta \leq \alpha$} \textbf{break} \EndIf
        \EndFor
    \EndIf
    \State \Return $v$
\EndFunction
\end{algorithmic}
\end{algorithm}


\subsection{Q-Learning (Tabular)}

Q-learning is a model-free, off-policy RL algorithm that learns an action-value function
$Q(s, a)$ estimating the expected cumulative reward from taking action $a$ in state $s$ and
following the optimal policy thereafter~\cite{watkins1992,sutton2018}. After each action, the
Q-value is updated:
\begin{equation}
    Q(s, a) \leftarrow Q(s, a) + \alpha \bigl[ r + \gamma \max_{a'} Q(s', a') - Q(s, a) \bigr]
\end{equation}
where $\alpha$ is the learning rate, $\gamma$ the discount factor, and $r$ the immediate reward.

Exploration uses $\varepsilon$-greedy: with probability $\varepsilon$ a random valid move is
selected; otherwise the highest-Q move is chosen. $\varepsilon$ decays exponentially from 1.0 to
0.05 over training. The state representation is the flattened board as a tuple (9 elements for
TTT, 42 for C4). This is exact for TTT but far too sparse for C4.

\begin{algorithm}[H]
\caption{Q-Learning Training Loop}
\begin{algorithmic}[1]
\State Initialise $Q(s,a) \leftarrow 0$ for all $s, a$
\For{episode $= 1, \ldots, N$}
    \State $s \leftarrow$ initial state
    \While{game not terminal}
        \State $a \leftarrow \varepsilon\text{-greedy}(Q, s)$
        \State Execute $a$, observe $r$, next state $s'$
        \State $Q(s,a) \leftarrow Q(s,a) + \alpha[r + \gamma \max_{a'} Q(s',a') - Q(s,a)]$
        \State $s \leftarrow s'$
    \EndWhile
    \State Decay $\varepsilon$
\EndFor
\end{algorithmic}
\end{algorithm}


\subsection{Deep Q-Network (DQN)}

DQN replaces the Q-table with a neural network $Q_\theta(s, a)$, enabling generalisation
across unseen states~\cite{mnih2015}. This is critical for Connect~4, where tabular coverage
is impractical.

The network takes the flattened board as input (9 for TTT, 42 for C4) and outputs Q-values for
each action (9 or 7). Two hidden layers with ReLU activations are used (default: 128, 64
neurons). Two techniques stabilise training:

\begin{itemize}
    \item \textbf{Experience replay}: transitions $(s, a, r, s', \text{done})$ are stored in a
          buffer. Training batches are sampled uniformly, breaking temporal correlations.
    \item \textbf{Target network}: a separate copy updated every $C$ steps provides stable
          targets:
          \begin{equation}
              \mathcal{L}(\theta) = \mathbb{E}\bigl[(r + \gamma \max_{a'} Q_{\theta^-}(s', a')
              - Q_\theta(s, a))^2\bigr]
          \end{equation}
\end{itemize}

Invalid moves are masked during action selection: only Q-values for valid moves are considered.


\subsection{Connect 4 Evaluation Function}
\label{sec:eval_fn}

Since full-depth search is infeasible for Connect~4 (Section~\ref{sec:scalability}), Minimax
uses a depth limit with a heuristic evaluation. The function scores every window of four
consecutive cells (horizontal, vertical, both diagonals):

\begin{table}[H]
\centering
\caption{Heuristic window scores for the Connect 4 evaluation function.}
\label{tab:eval_scores}
\small
\begin{tabular}{lc}
\toprule
\textbf{Window configuration} & \textbf{Score} \\
\midrule
4 own pieces (win)                   & $+100$ \\
3 own + 1 empty                      & $+5$ \\
2 own + 2 empty                      & $+2$ \\
3 opponent + 1 empty (threat)        & $-4$ \\
\bottomrule
\end{tabular}
\end{table}

Centre-column pieces receive a $+3$ bonus each, reflecting the centre's strategic
importance~\cite{allis1988}. The function sums scores across all windows and returns the total.


% ══════════════════════════════════════════════════════════════════════════════
\section{Experimental Setup}

\paragraph{Platform.}
All experiments ran in Google Colab (Python~3.10, PyTorch). DQN training used a T4 GPU. All
code is in a single script (\texttt{cs7is2\_assignment3.py}) with a unified entry point.

\paragraph{Games and agents.}
\begin{itemize}
    \item \textbf{Tic Tac Toe}: Minimax, Minimax+AB, Q-Learning, DQN, Default Opponent.
    \item \textbf{Connect 4}: Minimax+AB (depth-limited), Q-Learning, DQN, Default Opponent.
\end{itemize}

\paragraph{Training protocol.}
For TTT, Q-Learning (50k episodes) and DQN (30k episodes) train against the default opponent.
For C4, as recommended by the specification, RL agents train against a \textbf{random} opponent
(Q-Learning: 100k episodes; DQN: 50k episodes).

\paragraph{Evaluation protocol.}
200 games per TTT evaluation vs default; 100 games per round-robin pair (TTT); 50 games per
C4 matchup. Player~1 is consistent across each matchup.

\paragraph{Metrics.}
\begin{itemize}
    \item \textbf{Win/draw/loss rate}: fraction of games won, drawn, lost.
    \item \textbf{Nodes visited}: total nodes expanded during minimax search.
    \item \textbf{Average move time}: wall-clock time per move.
    \item \textbf{Training convergence}: rolling win rate over episodes for RL agents.
    \item \textbf{Parameter sensitivity}: performance under different hyperparameters.
\end{itemize}


% ══════════════════════════════════════════════════════════════════════════════
\section{Connect 4: Scalability Analysis}
\label{sec:scalability}

\subsection{Full-Depth Infeasibility}

To confirm full-depth Minimax is infeasible for Connect~4, both variants were executed on an
empty board and timed.

% ┌─────────────────────────────────────────────────────────────────────┐
% │  FILL IN: Replace X, Y with your actual node counts from the      │
% │  scalability test output (Phase 1, c4_scalability_test)            │
% └─────────────────────────────────────────────────────────────────────┘

\begin{table}[H]
\centering
\caption{Full-depth Minimax on an empty Connect 4 board (terminated after 30 minutes).}
\label{tab:scalability}
\small
\begin{tabular}{lrrr}
\toprule
\textbf{Algorithm} & \textbf{Time (s)} & \textbf{Nodes visited} & \textbf{Moves completed} \\
\midrule
Minimax (no pruning)   & 1800 & \textit{X} & 0 of 42 \\
Minimax (alpha-beta)   & 1800 & \textit{Y} & 0 of 42 \\
\bottomrule
\end{tabular}
\end{table}

Neither variant completed even the first move within 30 minutes. The C4 game tree has branching
factor~7 and maximum depth~42; the total leaf count is approximately $7^{42} \approx
3.1 \times 10^{35}$. Alpha-beta pruning reduces nodes visited but not enough for full-depth
feasibility.


\subsection{Depth-Limited Search with Evaluation Function}

To make Minimax tractable for C4, search is terminated at a depth limit with the heuristic
evaluation function from Section~\ref{sec:eval_fn}. Depths 2--7 were tested, each playing
20 games against the default opponent.

\begin{figure}[H]
    \centering
    \tryimg[width=\linewidth]{c4_depth_analysis}
    \caption{Connect 4 depth-limited Minimax: win rate (left), avg move time (centre), and
    avg nodes visited (right) at depth limits 2--7.}
    \label{fig:depth_analysis}
\end{figure}

% ┌─────────────────────────────────────────────────────────────────────┐
% │  FILL IN: Replace -- with actual values from Phase 8 output        │
% └─────────────────────────────────────────────────────────────────────┘

\begin{table}[H]
\centering
\caption{Connect 4 Minimax+AB performance at different depth limits vs default opponent.}
\label{tab:depth}
\small
\begin{tabular}{rrrr}
\toprule
\textbf{Depth} & \textbf{Win Rate (\%)} & \textbf{Avg Time (ms)} & \textbf{Avg Nodes} \\
\midrule
2 & -- & -- & -- \\
3 & -- & -- & -- \\
4 & -- & -- & -- \\
5 & -- & -- & -- \\
6 & -- & -- & -- \\
7 & -- & -- & -- \\
\bottomrule
\end{tabular}
\end{table}

Depth~5 was selected for subsequent experiments, offering strong play with reasonable
computation time. Beyond depth~5, time increases substantially while win-rate gains diminish.


\subsection{RL Training Considerations for Connect 4}

The C4 state space is too large for tabular Q-learning to visit meaningfully. A 100k-episode
run encounters at most $\sim$2.1M states ($\sim$21 moves $\times$ 100k), a vanishing fraction
of $\sim$4.5 trillion total. As recommended by the specification, C4 RL agents train against
a \textbf{random opponent}, which is easier to beat and allows even partially trained agents
to show improvement.

DQN is expected to outperform tabular Q-learning on C4 because the neural network generalises
across structurally similar board positions.


% ══════════════════════════════════════════════════════════════════════════════
\section{Results: Tic Tac Toe}

\subsection{Algorithms vs Default Opponent}

% ┌─────────────────────────────────────────────────────────────────────┐
% │  FILL IN: W/D/L from ttt_vs_default dict (Phase 3 output)         │
% └─────────────────────────────────────────────────────────────────────┘

\begin{table}[H]
\centering
\caption{TTT: each algorithm (player 1) vs Default Opponent (player $-1$) over 200 games.}
\label{tab:ttt_vs_default}
\small
\begin{tabular}{lrrrr}
\toprule
\textbf{Algorithm} & \textbf{Wins} & \textbf{Draws} & \textbf{Losses} & \textbf{Win Rate} \\
\midrule
Minimax              & -- & -- & -- & -- \\
Minimax + AB         & -- & -- & -- & -- \\
Q-Learning           & -- & -- & -- & -- \\
DQN                  & -- & -- & -- & -- \\
\bottomrule
\end{tabular}
\end{table}

Both Minimax variants achieve identical results: alpha-beta does not change the selected move,
only the efficiency of finding it. Minimax on TTT with full depth is provably optimal; it never
loses~\cite{russell2010}. Any non-wins are draws, since TTT with perfect play from both sides
always draws.

Q-Learning and DQN performance depends on training quality. After 50k and 30k episodes
respectively, both should approach near-optimal play, though minor imperfections may persist
from insufficiently explored states.


\subsection{Head-to-Head Round Robin}

A round-robin tournament was run with all four algorithms plus the default opponent. Each pair
played 100 games with the row agent as player~1.

\begin{figure}[H]
    \centering
    \tryimg[width=0.75\linewidth]{ttt_round_robin}
    \caption{TTT round-robin: win rate of row player (player 1) vs column player. Diagonal = 50\%.}
    \label{fig:ttt_roundrobin}
\end{figure}

% ┌─────────────────────────────────────────────────────────────────────┐
% │  FILL IN: 2-3 sentences interpreting the heatmap.                  │
% │  e.g. "Minimax agents dominate..." or "Q-Learning and DQN show    │
% │  comparable performance..."                                        │
% └─────────────────────────────────────────────────────────────────────┘


\subsection{Alpha-Beta Pruning Efficiency}

Both Minimax variants were run on an empty TTT board (worst case: all 9 cells available).

% ┌─────────────────────────────────────────────────────────────────────┐
% │  FILL IN: values from Phase 9 console output                       │
% └─────────────────────────────────────────────────────────────────────┘

\begin{table}[H]
\centering
\caption{Nodes visited on an empty Tic Tac Toe board.}
\label{tab:ab_savings}
\small
\begin{tabular}{lrr}
\toprule
\textbf{Algorithm} & \textbf{Nodes visited} & \textbf{Reduction} \\
\midrule
Minimax (no pruning) & -- & — \\
Minimax (alpha-beta) & -- & --\% \\
\bottomrule
\end{tabular}
\end{table}

Alpha-beta achieves a substantial reduction without changing the selected move. This saving is
even more significant for Connect~4, where it determines whether depth-5+ search is practical.


\subsection{RL Training Curves}

\begin{figure}[H]
    \centering
    \begin{subfigure}[b]{0.48\linewidth}
        \tryimg[width=\linewidth]{ttt_ql_training}
        \caption{Q-Learning (50k episodes)}
    \end{subfigure}
    \hfill
    \begin{subfigure}[b]{0.48\linewidth}
        \tryimg[width=\linewidth]{ttt_dqn_training}
        \caption{DQN (30k episodes)}
    \end{subfigure}
    \caption{TTT RL training curves: win, loss, and draw rates (rolling avg, window=500).}
    \label{fig:ttt_training}
\end{figure}

% ┌─────────────────────────────────────────────────────────────────────┐
% │  FILL IN: 2-3 sentences on convergence. When does win rate         │
% │  stabilise? Does loss rate drop to ~0? Q-Learning vs DQN speed?    │
% └─────────────────────────────────────────────────────────────────────┘


% ══════════════════════════════════════════════════════════════════════════════
\section{Results: Connect 4}

\subsection{Algorithms vs Default Opponent}

% ┌─────────────────────────────────────────────────────────────────────┐
% │  FILL IN: values from c4_vs_default dict (Phase 6 output)          │
% └─────────────────────────────────────────────────────────────────────┘

\begin{table}[H]
\centering
\caption{C4: each algorithm (player 1) vs Default Opponent (player $-1$) over 50 games.}
\label{tab:c4_vs_default}
\small
\begin{tabular}{lrrrr}
\toprule
\textbf{Algorithm} & \textbf{Wins} & \textbf{Draws} & \textbf{Losses} & \textbf{Win Rate} \\
\midrule
Minimax+AB (d=5) & -- & -- & -- & -- \\
Q-Learning       & -- & -- & -- & -- \\
DQN              & -- & -- & -- & -- \\
\bottomrule
\end{tabular}
\end{table}

% ┌─────────────────────────────────────────────────────────────────────┐
% │  FILL IN: analysis — does minimax dominate? How does Q-Learning    │
% │  compare to DQN? Do RL agents trained vs random struggle vs        │
% │  the smarter default opponent?                                     │
% └─────────────────────────────────────────────────────────────────────┘


\subsection{Head-to-Head Round Robin}

\begin{figure}[H]
    \centering
    \tryimg[width=0.7\linewidth]{c4_round_robin}
    \caption{C4 round-robin: win rate of row player vs column player over 50 games per pair.}
    \label{fig:c4_roundrobin}
\end{figure}

% ┌─────────────────────────────────────────────────────────────────────┐
% │  FILL IN: 2-3 sentences interpreting the C4 heatmap                │
% └─────────────────────────────────────────────────────────────────────┘


\subsection{RL Training Curves}

\begin{figure}[H]
    \centering
    \begin{subfigure}[b]{0.48\linewidth}
        \tryimg[width=\linewidth]{c4_ql_training}
        \caption{Q-Learning (100k episodes)}
    \end{subfigure}
    \hfill
    \begin{subfigure}[b]{0.48\linewidth}
        \tryimg[width=\linewidth]{c4_dqn_training}
        \caption{DQN (50k episodes)}
    \end{subfigure}
    \caption{C4 RL training curves (vs random opponent, rolling avg window=1000).}
    \label{fig:c4_training}
\end{figure}

% ┌─────────────────────────────────────────────────────────────────────┐
% │  FILL IN: convergence speed, final win rate, Q-Learning vs DQN     │
% └─────────────────────────────────────────────────────────────────────┘


% ══════════════════════════════════════════════════════════════════════════════
\section{Parameter Experimentation}
\label{sec:param_exp}

A key critique of the previous assignment was the absence of parameter sensitivity analysis.
This section addresses that by sweeping key hyperparameters for each RL algorithm.


\subsection{Q-Learning Hyperparameters}

Six configurations were tested on TTT (20k episodes each), varying learning rate $\alpha$,
discount factor $\gamma$, and epsilon decay rate:

\begin{table}[H]
\centering
\caption{Q-Learning parameter configurations tested on Tic Tac Toe.}
\label{tab:ql_params}
\small
\begin{tabular}{llll}
\toprule
\textbf{Label} & $\boldsymbol{\alpha}$ & $\boldsymbol{\gamma}$ & \textbf{$\varepsilon$ decay} \\
\midrule
$\alpha$=0.05, $\gamma$=0.9 & 0.05 & 0.9 & 0.999 \\
$\alpha$=0.1, $\gamma$=0.9  & 0.1  & 0.9 & 0.999 \\
$\alpha$=0.2, $\gamma$=0.9  & 0.2  & 0.9 & 0.999 \\
$\alpha$=0.1, $\gamma$=0.5  & 0.1  & 0.5 & 0.999 \\
$\alpha$=0.1, $\gamma$=0.99 & 0.1  & 0.99 & 0.999 \\
$\alpha$=0.1, slow decay    & 0.1  & 0.9  & 0.9999 \\
\bottomrule
\end{tabular}
\end{table}

\begin{figure}[H]
    \centering
    \tryimg[width=\linewidth]{ttt_ql_param_sweep}
    \caption{Q-Learning parameter sensitivity on TTT. Each line shows rolling win rate
    (window=500) for a different hyperparameter configuration.}
    \label{fig:ql_params}
\end{figure}

% ┌─────────────────────────────────────────────────────────────────────┐
% │  FILL IN after viewing figure:                                      │
% │  - Which alpha works best? (0.1 balances speed + stability?)       │
% │  - Impact of gamma: 0.5 undervalues future? 0.99 slows learning?  │
% │  - Slow decay: explores longer — better or worse?                  │
% └─────────────────────────────────────────────────────────────────────┘


\subsection{DQN Architecture and Learning Rate}

Four DQN configurations were tested on TTT (15k episodes each):

\begin{table}[H]
\centering
\caption{DQN parameter configurations tested on Tic Tac Toe.}
\label{tab:dqn_params}
\small
\begin{tabular}{llll}
\toprule
\textbf{Label} & \textbf{Hidden layers} & \textbf{Learning rate} & \textbf{$\sim$Params} \\
\midrule
lr=1e-3, small  & 64, 32   & $10^{-3}$ & 2.8k \\
lr=1e-3, medium & 128, 64  & $10^{-3}$ & 10k \\
lr=1e-4, medium & 128, 64  & $10^{-4}$ & 10k \\
lr=1e-3, large  & 256, 128 & $10^{-3}$ & 37k \\
\bottomrule
\end{tabular}
\end{table}

\begin{figure}[H]
    \centering
    \tryimg[width=\linewidth]{ttt_dqn_param_sweep}
    \caption{DQN parameter sensitivity on TTT. Each line shows rolling win rate (window=500).}
    \label{fig:dqn_params}
\end{figure}

% ┌─────────────────────────────────────────────────────────────────────┐
% │  FILL IN after viewing figure:                                      │
% │  - Larger network help or hurt for small TTT game?                 │
% │  - lr 1e-3 vs 1e-4: convergence speed vs stability?                │
% └─────────────────────────────────────────────────────────────────────┘


\subsection{Minimax Depth Limit (Connect 4)}

The depth analysis was presented in Section~\ref{sec:scalability}
(Figure~\ref{fig:depth_analysis}, Table~\ref{tab:depth}). Depth~5 was selected as the
operating point: strong play against the default opponent with sub-second move times.


% ══════════════════════════════════════════════════════════════════════════════
\section{Overall Comparison}

This section synthesises findings across both games. The three algorithm families differ
fundamentally in their trade-offs.

\begin{table}[H]
\centering
\caption{Summary of algorithm properties across both games.}
\label{tab:summary}
\small
\begin{tabular}{p{2.6cm}p{3.5cm}p{3.5cm}p{3.5cm}}
\toprule
\textbf{Property} & \textbf{Minimax (+AB)} & \textbf{Q-Learning} & \textbf{DQN} \\
\midrule
Optimality &
    Provably optimal at full depth &
    Converges to optimal given infinite episodes &
    Approximates optimal via function approximation \\
\addlinespace
TTT performance &
    Perfect play (never loses) &
    Near-perfect after 50k episodes &
    Near-perfect after 30k episodes \\
\addlinespace
C4 feasibility &
    Depth-limited only &
    State space too large for coverage &
    Generalises across unseen states \\
\addlinespace
Training cost &
    None &
    50k--100k episodes, CPU minutes &
    30k--50k episodes, GPU helps \\
\addlinespace
Move time &
    Depth-dependent (ms--seconds) &
    $O(1)$ table lookup &
    $O(1)$ forward pass \\
\addlinespace
Memory &
    $O(b^d)$ stack frames &
    $O(|S| \times |A|)$ table &
    Fixed network parameters \\
\bottomrule
\end{tabular}
\end{table}

\paragraph{Tic Tac Toe.}
Minimax is the clear winner: provably optimal play with negligible computation and zero
training cost. Both RL agents approach comparable performance after training but require
thousands of episodes to learn what Minimax computes directly from the game tree. The TTT
Q-table is small enough ($\leq$5,478 states $\times$ 9 actions) that tabular Q-learning is
adequate; DQN provides no meaningful advantage.

\paragraph{Connect 4.}
The picture inverts. Full-depth Minimax is infeasible; depth-limited search at depth~5
requires evaluating thousands of nodes per move. Q-Learning struggles because the state space
is too large for meaningful coverage. DQN generalises best across unseen states. However,
depth-limited Minimax with a good evaluation function still provides the strongest play per
move, as it explicitly considers the adversarial structure.

\paragraph{Paradigm trade-off.}
The fundamental distinction is between \textit{compute at play time} (Minimax) and
\textit{compute at training time} (RL). Minimax requires no preparation but searches the tree
per move. RL invests computation upfront and then decides in constant time. For tractable games
(TTT), Minimax is superior. For larger games (C4), RL that generalises is necessary, though
depth-limited search with heuristics remains competitive.


% ══════════════════════════════════════════════════════════════════════════════
\section{Conclusions}

% ┌─────────────────────────────────────────────────────────────────────┐
% │  FILL IN: Replace placeholder numbers (X, Y, Z, W, T) with your   │
% │  actual experimental results after running the code.               │
% └─────────────────────────────────────────────────────────────────────┘

\begin{enumerate}
    \item \textbf{Minimax achieves perfect play on Tic Tac Toe.} Both pruned and unpruned
          variants never lose, confirming full-depth Minimax is optimal for tractable games.

    \item \textbf{Alpha-beta pruning dramatically reduces computation.} On an empty TTT board,
          pruning reduces nodes from \textit{X} to \textit{Y} (a \textit{Z}\% reduction) while
          selecting the identical move.

    \item \textbf{Full-depth Minimax is infeasible for Connect 4.} After 30 minutes, neither
          variant completed the first move. The branching factor of 7 and depth of 42 produce
          $\sim$$3.1 \times 10^{35}$ leaf nodes.

    \item \textbf{Depth-limited Minimax at depth 5 provides strong C4 play.} With alpha-beta
          and a heuristic evaluation function, this achieves \textit{W}\% win rate vs default
          with \textit{T} ms average move time.

    \item \textbf{Q-Learning converges on TTT but struggles on C4.} Tabular RL reaches
          near-optimal TTT play after $\sim$50k episodes, but the C4 state space is too
          large for meaningful Q-table coverage.

    \item \textbf{DQN generalises better than tabular Q-Learning on C4.} The neural network
          shares learned features across similar board states, allowing reasonable play in
          unseen positions.

    \item \textbf{Hyperparameter choice significantly impacts RL.} $\alpha = 0.1$,
          $\gamma = 0.9$ provided the best Q-Learning balance. For DQN, $\text{lr} = 10^{-3}$
          with a medium network (128, 64) converged fastest.

    \item \textbf{The optimal algorithm depends on game complexity.} For small, tractable
          games, Minimax is superior. For larger games, depth-limited Minimax and DQN
          represent the most practical approaches.
\end{enumerate}


% ══════════════════════════════════════════════════════════════════════════════
\newpage
\begin{thebibliography}{10}

\bibitem{vonneumann1944}
J.~von~Neumann and O.~Morgenstern, \textit{Theory of Games and Economic Behavior}.
Princeton University Press, 1944.

\bibitem{russell2010}
S.~J.~Russell and P.~Norvig, \textit{Artificial Intelligence: A Modern Approach},
3rd ed. Prentice Hall, 2010.

\bibitem{knuth1975}
D.~E.~Knuth and R.~W.~Moore,
``An analysis of alpha-beta pruning,''
\textit{Artificial Intelligence}, vol.~6, no.~4, pp.~293--326, 1975.

\bibitem{watkins1992}
C.~J.~C.~H.~Watkins and P.~Dayan,
``Q-learning,''
\textit{Machine Learning}, vol.~8, no.~3, pp.~279--292, 1992.

\bibitem{mnih2015}
V.~Mnih \textit{et al.},
``Human-level control through deep reinforcement learning,''
\textit{Nature}, vol.~518, pp.~529--533, 2015.

\bibitem{sutton2018}
R.~S.~Sutton and A.~G.~Barto, \textit{Reinforcement Learning: An Introduction},
2nd ed. MIT Press, 2018.

\bibitem{allis1988}
L.~V.~Allis,
``A knowledge-based approach of Connect-Four,''
Master's thesis, Vrije Universiteit Amsterdam, 1988.

\bibitem{bellman1957}
R.~Bellman, ``A Markovian decision process,''
\textit{Journal of Mathematics and Mechanics}, vol.~6, no.~5, pp.~679--684, 1957.

\end{thebibliography}


% ══════════════════════════════════════════════════════════════════════════════
\newpage
\appendix
\section{Appendix: Algorithm Implementations}
\label{sec:appendix}

The specification requires code for six algorithm implementations (2 games $\times$ 3
algorithms) as appendices.

\subsection{A1: Minimax --- Tic Tac Toe}
\begin{lstlisting}[language=Python, caption={Minimax (no pruning) applied to Tic Tac Toe.}]
class MinimaxAgent:
    name = "Minimax"

    def __init__(self, player=1, depth_limit=None):
        self.player = player
        self.depth_limit = depth_limit
        self.nodes_visited = 0

    def choose_move(self, game, training=False):
        self.nodes_visited = 0
        valid = game.get_valid_actions()
        best_action = valid[0]
        best_val = -math.inf
        for action in valid:
            game.make_move(action)
            val = self._minimax(game, depth=1, maximizing=False)
            game.undo_move(action)
            if val > best_val:
                best_val = val
                best_action = action
        return best_action

    def _minimax(self, game, depth, maximizing):
        self.nodes_visited += 1
        winner = game.get_winner()
        if winner is not None:
            return self._terminal_value(winner)
        if self.depth_limit is not None and depth >= self.depth_limit:
            return self._heuristic(game)
        valid = game.get_valid_actions()
        if maximizing:
            best = -math.inf
            for action in valid:
                game.make_move(action)
                best = max(best, self._minimax(game, depth+1, False))
                game.undo_move(action)
            return best
        else:
            best = math.inf
            for action in valid:
                game.make_move(action)
                best = min(best, self._minimax(game, depth+1, True))
                game.undo_move(action)
            return best

    def _terminal_value(self, winner):
        if winner == self.player:   return  1.0
        elif winner == 0:           return  0.0
        else:                       return -1.0
\end{lstlisting}

\subsection{A2: Minimax with Alpha-Beta --- Connect 4}
\begin{lstlisting}[language=Python, caption={Minimax+AB (depth-limited) applied to Connect 4.}]
class MinimaxABAgent:
    name = "Minimax+AB"

    def __init__(self, player=1, depth_limit=None):
        self.player = player
        self.depth_limit = depth_limit
        self.nodes_visited = 0

    def choose_move(self, game, training=False):
        self.nodes_visited = 0
        valid = self._order_moves(game, game.get_valid_actions())
        best_action = valid[0]
        best_val = -math.inf
        alpha, beta = -math.inf, math.inf
        for action in valid:
            game.make_move(action)
            val = self._alphabeta(game, 1, alpha, beta, False)
            game.undo_move(action)
            if val > best_val:
                best_val = val
                best_action = action
            alpha = max(alpha, best_val)
        return best_action

    def _alphabeta(self, game, depth, alpha, beta, maximizing):
        self.nodes_visited += 1
        winner = game.get_winner()
        if winner is not None:
            return self._terminal_value(winner)
        if self.depth_limit is not None and depth >= self.depth_limit:
            return self._heuristic(game)
        valid = self._order_moves(game, game.get_valid_actions())
        if maximizing:
            value = -math.inf
            for action in valid:
                game.make_move(action)
                value = max(value,
                    self._alphabeta(game, depth+1, alpha, beta, False))
                game.undo_move(action)
                alpha = max(alpha, value)
                if alpha >= beta:
                    break   # beta cut-off
            return value
        else:
            value = math.inf
            for action in valid:
                game.make_move(action)
                value = min(value,
                    self._alphabeta(game, depth+1, alpha, beta, True))
                game.undo_move(action)
                beta = min(beta, value)
                if beta <= alpha:
                    break   # alpha cut-off
            return value

    def _order_moves(self, game, valid):
        # Centre-first ordering for Connect 4
        if isinstance(game, Connect4):
            ordered = [c for c in Connect4.CENTER_ORDER if c in valid]
            ordered += [c for c in valid if c not in ordered]
            return ordered
        return valid

    def _heuristic(self, game):
        if isinstance(game, Connect4):
            return game.evaluate_heuristic(self.player) / 200.0
        return 0.0
\end{lstlisting}

\subsection{A3: Q-Learning --- Tic Tac Toe}
\begin{lstlisting}[language=Python, caption={Tabular Q-Learning agent and TTT training configuration.}]
class QLearningAgent:
    name = "Q-Learning"

    def __init__(self, player=1, alpha=0.1, gamma=0.99,
                 epsilon=1.0, epsilon_min=0.05, epsilon_decay=0.9995):
        self.player = player
        self.alpha, self.gamma = alpha, gamma
        self.epsilon, self.epsilon_min = epsilon, epsilon_min
        self.epsilon_decay = epsilon_decay
        self.q_table = defaultdict(lambda: defaultdict(float))

    def choose_move(self, game, training=False):
        state = game.get_state_key(perspective=self.player)
        valid = game.get_valid_actions()
        if training and random.random() < self.epsilon:
            return random.choice(valid)
        best_val, best_action = -float("inf"), None
        for action in valid:
            q = self.q_table[state][action]
            if q > best_val:
                best_val, best_action = q, action
        return best_action if best_action is not None else random.choice(valid)

    def update(self, state, action, reward, next_state, next_valid, done):
        current_q = self.q_table[state][action]
        if done or not next_valid:
            target = reward
        else:
            future_qs = [self.q_table[next_state][a] for a in next_valid]
            target = reward + self.gamma * max(future_qs)
        self.q_table[state][action] = (current_q
            + self.alpha * (target - current_q))

    def decay_epsilon(self):
        self.epsilon = max(self.epsilon_min,
                           self.epsilon * self.epsilon_decay)

# TTT training configuration
ttt_ql = QLearningAgent(player=1, alpha=0.1, gamma=0.99,
                        epsilon=1.0, epsilon_min=0.05,
                        epsilon_decay=0.9995)
trainer = Trainer(TicTacToe, ttt_ql, DefaultAgent(),
                  n_episodes=50_000, swap_sides=True)
metrics = trainer.train()
\end{lstlisting}

\subsection{A4: Q-Learning --- Connect 4}
\begin{lstlisting}[language=Python, caption={Q-Learning Connect 4 training configuration (vs random opponent).}]
# C4 Q-Learning training configuration
# Trained vs Random (easier target; state space too large for Default)
c4_ql = QLearningAgent(player=1, alpha=0.1, gamma=0.99,
                       epsilon=1.0, epsilon_min=0.05,
                       epsilon_decay=0.9999)
trainer = Trainer(Connect4, c4_ql, RandomAgent(),
                  n_episodes=100_000, swap_sides=True)
metrics = trainer.train()

# Note: Q-table covers only a tiny fraction of the ~4.5e12 C4 states.
# A 100k-episode run visits at most ~2.1M states. The agent learns
# limited local patterns but cannot generalise to unseen positions.
\end{lstlisting}

\subsection{A5: DQN --- Tic Tac Toe}
\begin{lstlisting}[language=Python, caption={DQN network, agent, and TTT training configuration.}]
class DQNNetwork(nn.Module):
    def __init__(self, input_size, hidden_sizes, output_size):
        super().__init__()
        layers = []
        prev = input_size
        for h in hidden_sizes:
            layers += [nn.Linear(prev, h), nn.ReLU()]
            prev = h
        layers.append(nn.Linear(prev, output_size))
        self.net = nn.Sequential(*layers)

    def forward(self, x):
        return self.net(x)

class DQNAgent:
    name = "DQN"

    def __init__(self, player=1, state_size=9, action_size=9,
                 hidden_sizes=(128, 128), lr=1e-3, gamma=0.99,
                 epsilon=1.0, epsilon_min=0.05, epsilon_decay=0.995,
                 batch_size=64, memory_size=50_000,
                 target_update_freq=200, game_name="tictactoe"):
        self.player = player
        self.gamma, self.epsilon = gamma, epsilon
        self.epsilon_min, self.epsilon_decay = epsilon_min, epsilon_decay
        self.batch_size, self.steps = batch_size, 0
        self.target_update_freq = target_update_freq
        self.device = torch.device("cpu")  # TTT: always CPU
        self.policy_net = DQNNetwork(
            state_size, hidden_sizes, action_size).to(self.device)
        self.target_net = DQNNetwork(
            state_size, hidden_sizes, action_size).to(self.device)
        self.target_net.load_state_dict(self.policy_net.state_dict())
        self.optimizer = optim.Adam(
            self.policy_net.parameters(), lr=lr)
        self.memory = deque(maxlen=memory_size)

    def choose_move(self, game, training=False):
        valid = game.get_valid_actions()
        if training and random.random() < self.epsilon:
            return random.choice(valid)
        arr = game.get_state_array(perspective=self.player)
        state = torch.FloatTensor(arr).unsqueeze(0).to(self.device)
        with torch.no_grad():
            q_values = self.policy_net(state).squeeze().cpu()
        masked = torch.full((self.policy_net.net[-1].out_features,),
                            float("-inf"))
        for a in valid:
            masked[a] = q_values[a]
        return int(masked.argmax().item())

# TTT DQN training configuration
ttt_dqn = DQNAgent(player=1, state_size=9, action_size=9,
                   hidden_sizes=(128, 128), lr=1e-3, gamma=0.99,
                   epsilon=1.0, epsilon_min=0.05,
                   epsilon_decay=0.995, batch_size=64,
                   memory_size=20_000, target_update_freq=200,
                   game_name="tictactoe")
trainer = Trainer(TicTacToe, ttt_dqn, DefaultAgent(),
                  n_episodes=30_000, swap_sides=True, train_every=1)
metrics = trainer.train()
\end{lstlisting}

\subsection{A6: DQN --- Connect 4}
\begin{lstlisting}[language=Python, caption={DQN Connect 4 training configuration (GPU if available).}]
# C4 DQN training configuration
# Larger network + GPU for Connect 4's 42-dimensional state space
c4_dqn = DQNAgent(player=1, state_size=42, action_size=7,
                  hidden_sizes=(256, 128, 64), lr=5e-4, gamma=0.99,
                  epsilon=1.0, epsilon_min=0.05,
                  epsilon_decay=0.9998, batch_size=128,
                  memory_size=50_000, target_update_freq=500,
                  game_name="connect4")
# game_name="connect4" triggers GPU use if torch.cuda.is_available()

trainer = Trainer(Connect4, c4_dqn, RandomAgent(),
                  n_episodes=50_000, swap_sides=True, train_every=1)
metrics = trainer.train()

# The neural network generalises across structurally similar C4 positions,
# overcoming the coverage limitation of tabular Q-learning.
\end{lstlisting}


\end{document}
